\section{Introduction}\label{sec:introduction}
A service agreement can return to the same public operating condition after different sequences of renewal reviews. Those histories need not permit the same future service. A customer who could have exited at an earlier review may have accepted a different remaining payment obligation. A provider that remembers only the current public condition can therefore lose value, even when demand is Markovian, prices are fixed, and neither party has private information. The operational question is not whether full history improves a particular schedule. It is how to construct an accepted optimal protocol, how much information its execution actually needs, and what an optimized memory restriction costs.

We answer these questions for positive additive payments, separable strictly concave quadratic operating rewards, and an exogenous recombining public graph. The principal algorithm excludes intertemporal switching costs and coupled vector commitments. These conditions make one accumulated payment price sufficient for optimal continuation. Each participation cap replaces prices below a threshold by that threshold. Execution retains the largest threshold encountered, together with a query-specific initial price. We construct the entire response family before choosing the initial price; the construction neither obtains an extensive-form optimum first nor enumerates a scenario tree.

The closest optimization connection is classical tree-constrained resource allocation. Mjelde's tree constraints and Tang's separable convex reductions are directly relevant, not merely analogies \citep{Mjelde1983,Tang1990}. After weighting each occurrence by its probability and discount factor, our full-history model belongs to that class. Moreover, memoizing the \emph{normalized} subtree response produces the same graph recursion. We make this equivalence explicit rather than present marginal allocation or memoization as new principles. Our results quantify what the quotient representation makes available: a graph-sized universe of response breakpoints, a polynomial rational-arithmetic implementation, and an exactly minimal execution alphabet for a fixed contractual query. These are separate statements about representation, computation, and implementation.

The paper develops three complementary results. First, an occurrence-to-vertex theorem identifies the precise information discarded by the quotient. A coefficient-event implementation then gives explicit arithmetic, storage, and rational bit bounds. Second, a backward equivalence relation merges price states exactly when they prescribe the same complete future actions. This relation characterizes the smallest writable alphabet for an optimal protocol and shows that its classes are ordered by incoming price. Strict concavity extends the exact-optimum lower bound to randomized controllers under the same accounting of retained information. Third, a renewal family admits a complete value--memory frontier with heterogeneous probabilities and operating costs. At a fixed mean payment, increasing outside-option dispersion strictly raises the loss from an insufficient memory alphabet, although arbitrarily small dispersion can already change the exact alphabet requirement.

An optimized public table supplies the economic comparator. At zero release it has a graph-sized convex formulation. At a fixed positive release, the price compiler supplies an exact value and a supporting inequality, including coincident physical and contractual bounds. We do not equate an exact oracle with finite termination of the continuous outer table-design problem.

The computational study separates implementation verification from comparison with other algorithms. It includes an independently implemented occurrence-indexed laminar reduction, a generic sparse convex optimizer on unfoldable instances, and a primal promise-grid method that already works on the public graph. Larger examples distinguish response-table growth from rational coefficient growth. No claim of superiority over all policy-graph or stochastic dual dynamic programming methods is inferred from the size of a hypothetical scenario tree. The broader accepted-adaptation results on switching, comparator-adjusted rents, continuous states, and certified deployment remain available in the retained theory volumes, with their statements and proofs unchanged; they are not used to enlarge the domain of the finite-price theorem.
